% GENERATED FILE. Edit canonical lessons, then run npm run generate:books.
% canonical-source-hash: fnv1a64:066a919948a59f52
% canonical-lessons: HI-C294-taftis, HI-C294-surag, HI-C294-jasus, HI-C294-sabut, HI-C294-khojna

\chapter{The Investigation}
\label{ch:hi-the-investigation}
\hlchaptermodality{294}

\begin{chapteropening}
\textbf{By the end of this chapter:} \emph{I can say what the police found and what will happen next.}

Complete the last lesson of chapter 294: I can say what the police found and what will happen next.
\end{chapteropening}

\section[\texorpdfstring{\hi{तफ़्तीश} (taftīś)}{तफ़्तीश (taftīś)}]{\hi{तफ़्तीश} (taftīś) — an investigation}

\label{lesson:HI-C294-taftis}



\begin{quote}
Before the new one: say the Hindi for a crook, then the Hindi for fake.
\end{quote}

\subsection*{You'll want to know: \hi{तफ़्तीश}}
\textbf{\hi{तफ़्तीश}} — \emph{taftīś} — \textquotedblleft{}an investigation\textquotedblright{}.

\begin{quote}
\hi{अगर} \hi{तफ़्तीश} \hi{ठीक} \hi{से} \hi{हुई} \hi{तो} \hi{सच} \hi{सामने} \hi{आएगा।}
\end{quote}

\begin{quote}
\emph{agar taftīś ṭhīk se huī to sac sāmne āegā.}
\end{quote}

\begin{quote}
If the investigation is done properly, the truth will come out.
\end{quote}

Feminine.

\begin{cousinweb}
\textbf{\hi{तफ़्तीश}} (\emph{taftīś}) is usually traced, through Persian, to Arabic \emph{taftīsh}, \textquotedblleft{}a searching, an inspection\textquotedblright{}. It is the word of police reports: \textbf{\hi{पुलिस} \hi{तफ़्तीश} \hi{कर} \hi{रही} \hi{है}}, the police are investigating.
\end{cousinweb}

\begin{grammarlens}[title={What you've built}]
One more word for telling what happened.
\end{grammarlens}

\subsection*{Your turn}
\begin{itemize}
\raggedright
  \item \emph{Say it:} \emph{taftīś}
  \item \emph{Say it:} \emph{taftīś}, once more
  \item \emph{Say it:} \emph{agar taftīś ṭhīk se huī to sac sāmne āegā.} — If the investigation is done properly, the truth will come out.
  \item \emph{Recall:} say the Hindi for a crook, then the Hindi for fake, then say \emph{taftīś} again
\end{itemize}

\begin{tcolorbox}[breakable,colback=teal!4,colframe=teal!35!black,title={Before you move on}]
What does \hi{तफ़्तीश} mean? (\textquotedblleft{}An investigation\textquotedblright{}.) Say it once more.
\end{tcolorbox}
\section[\texorpdfstring{\hi{सुराग़} (surāġ)}{सुराग़ (surāġ)}]{\hi{सुराग़} (surāġ) — a clue; a trace}

\label{lesson:HI-C294-surag}



\begin{quote}
Before the new one: say the Hindi for fake, then the Hindi for an investigation.
\end{quote}

\subsection*{You'll want to know: \hi{सुराग़}}
\textbf{\hi{सुराग़}} — \emph{surāġ} — \textquotedblleft{}a clue; a trace\textquotedblright{}.

\begin{quote}
\hi{पुलिस} \hi{को} \hi{दरवाज़े} \hi{पर} \hi{एक} \hi{सुराग़} \hi{मिला।}
\end{quote}

\begin{quote}
\emph{pulis ko darvāze par ek surāġ milā.}
\end{quote}

\begin{quote}
The police found a clue on the door.
\end{quote}

Masculine.

\begin{cousinweb}
is usually traced to Persian \emph{surāġ}, \textquotedblleft{}a trace, a track\textquotedblright{}. It is what you follow: one small \textbf{\hi{सुराग़}} can lead to the thief.
\end{cousinweb}

\begin{grammarlens}[title={What you've built}]
One more word for telling what happened.
\end{grammarlens}

\subsection*{Your turn}
\begin{itemize}
\raggedright
  \item \emph{Say it:} \emph{surāġ}
  \item \emph{Say it:} \emph{surāġ}, once more
  \item \emph{Say it:} \emph{pulis ko darvāze par ek surāġ milā.} — The police found a clue on the door.
  \item \emph{Recall:} say the Hindi for fake, then the Hindi for an investigation, then say \emph{surāġ} again
\end{itemize}

\begin{tcolorbox}[breakable,colback=teal!4,colframe=teal!35!black,title={Before you move on}]
What does \hi{सुराग़} mean? (\textquotedblleft{}A clue; a trace\textquotedblright{}.) Say it once more.
\end{tcolorbox}
\section[\texorpdfstring{\hi{जासूस} (jāsūs)}{जासूस (jāsūs)}]{\hi{जासूस} (jāsūs) — a detective; a spy}

\label{lesson:HI-C294-jasus}



\begin{quote}
Before the new one: say the Hindi for an investigation, then the Hindi for a clue.
\end{quote}

\subsection*{You'll want to know: \hi{जासूस}}
\textbf{\hi{जासूस}} — \emph{jāsūs} — \textquotedblleft{}a detective; a spy\textquotedblright{}.

\begin{quote}
\hi{बचपन} \hi{में} \hi{अरुण} \hi{जासूस} \hi{बनना} \hi{चाहता} \hi{था।}
\end{quote}

\begin{quote}
\emph{bacpan meṁ Aruṇ jāsūs bannā cāhtā thā.}
\end{quote}

\begin{quote}
As a child Arun wanted to be a detective.
\end{quote}

Masculine or feminine, by the person meant.

\begin{cousinweb}
\textbf{\hi{जासूस}} (\emph{jāsūs}) is usually traced, through Persian, to Arabic \emph{jāsūs}, \textquotedblleft{}a spy\textquotedblright{}. In stories it is the detective who works out who did it.
\end{cousinweb}

\begin{grammarlens}[title={What you've built}]
One more word for telling what happened.
\end{grammarlens}

\subsection*{Your turn}
\begin{itemize}
\raggedright
  \item \emph{Say it:} \emph{jāsūs}
  \item \emph{Say it:} \emph{jāsūs}, once more
  \item \emph{Say it:} \emph{bacpan meṁ Aruṇ jāsūs bannā cāhtā thā.} — As a child Arun wanted to be a detective.
  \item \emph{Recall:} say the Hindi for an investigation, then the Hindi for a clue, then say \emph{jāsūs} again
\end{itemize}

\begin{tcolorbox}[breakable,colback=teal!4,colframe=teal!35!black,title={Before you move on}]
What does \hi{जासूस} mean? (\textquotedblleft{}A detective; a spy\textquotedblright{}.) Say it once more.
\end{tcolorbox}
\section[\texorpdfstring{\hi{सबूत} (sabūt)}{सबूत (sabūt)}]{\hi{सबूत} (sabūt) — evidence; proof}

\label{lesson:HI-C294-sabut}



\begin{quote}
Before the new one: say the Hindi for a clue, then the Hindi for a detective.
\end{quote}

\subsection*{You'll want to know: \hi{सबूत}}
\textbf{\hi{सबूत}} — \emph{sabūt} — \textquotedblleft{}evidence; proof\textquotedblright{}.

\begin{quote}
\hi{पुलिस} \hi{को} \hi{सबूत} \hi{चाहिए}, \hi{नहीं} \hi{तो} \hi{चोर} \hi{छूट} \hi{जाएगा।}
\end{quote}

\begin{quote}
\emph{pulis ko sabūt cāhiye, nahīṃ to cor chūṭ jāegā.}
\end{quote}

\begin{quote}
The police need evidence, or the thief will go free.
\end{quote}

Masculine.

\begin{cousinweb}
\textbf{\hi{सबूत}} (\emph{sabūt}) is usually traced, through Persian, to Arabic \emph{thubūt}, \textquotedblleft{}firmness, being established\textquotedblright{}: what makes a claim stand. Hindi also spells it \textbf{\hi{सुबूत}}.
\end{cousinweb}

\begin{grammarlens}[title={What you've built}]
One more word for telling what happened.
\end{grammarlens}

\subsection*{Your turn}
\begin{itemize}
\raggedright
  \item \emph{Say it:} \emph{sabūt}
  \item \emph{Say it:} \emph{sabūt}, once more
  \item \emph{Say it:} \emph{pulis ko sabūt cāhiye, nahīṃ to cor chūṭ jāegā.} — The police need evidence, or the thief will go free.
  \item \emph{Recall:} say the Hindi for a clue, then the Hindi for a detective, then say \emph{sabūt} again
\end{itemize}

\begin{tcolorbox}[breakable,colback=teal!4,colframe=teal!35!black,title={Before you move on}]
What does \hi{सबूत} mean? (\textquotedblleft{}Evidence; proof\textquotedblright{}.) Say it once more.
\end{tcolorbox}
\section[\texorpdfstring{\hi{खोजना} (khojnā)}{खोजना (khojnā)}]{\hi{खोजना} (khojnā) — to search for}

\label{lesson:HI-C294-khojna}



\begin{quote}
Before the new one: say the Hindi for a detective, then the Hindi for evidence.
\end{quote}

\subsection*{You'll want to know: \hi{खोजना}}
\textbf{\hi{खोजना}} — \emph{khojnā} — \textquotedblleft{}to search for\textquotedblright{}.

\begin{quote}
\hi{पुलिस} \hi{ने} \hi{पूरे} \hi{गाँव} \hi{में} \hi{चोर} \hi{को} \hi{खोजा।}
\end{quote}

\begin{quote}
\emph{pulis ne pūre gā̃v meṁ cor ko khojā.}
\end{quote}

\begin{quote}
The police searched the whole village for the thief.
\end{quote}

Past, with \hi{ने} (he or she …): \hi{उसने} \hi{चाबी} \hi{खोजी}, \emph{usne cābī khojī}.

\begin{cousinweb}
is inherited from \emph{*khōjja}, a form CDIAL reconstructs (CDIAL 3929), with the sense of following a track. It is close to \textbf{\hi{ढूँढना}}, to look for; \textbf{\hi{खोजना}} is the slower, more thorough search.
\end{cousinweb}

\begin{grammarlens}[title={What you've built}]
One more word for telling what happened.
\end{grammarlens}

\subsection*{Your turn}
\begin{itemize}
\raggedright
  \item \emph{Say it:} \emph{khojnā}
  \item \emph{Say it:} \emph{khojnā}, once more
  \item \emph{Say it:} \emph{pulis ne pūre gā̃v meṁ cor ko khojā.} — The police searched the whole village for the thief.
  \item \emph{Recall:} say the Hindi for a detective, then the Hindi for evidence, then say \emph{khojnā} again
\end{itemize}

\begin{tcolorbox}[breakable,colback=teal!4,colframe=teal!35!black,title={Before you move on}]
What does \hi{खोजना} mean? (\textquotedblleft{}To search for\textquotedblright{}.) Say it once more.
\end{tcolorbox}
